\documentclass[11pt]{article}

\usepackage{amsmath,amssymb}
\usepackage{xurl}
\usepackage[colorlinks=true,linkcolor=blue,citecolor=blue,urlcolor=blue]{hyperref}

\input{./command.tex}

\title{Wolf Theorem 6.3: Positive-Map Scope and Boundary Corrections}
\author{}
\date{2026-08-24}

\begin{document}
\maketitle
\gapnote{false-source}{resolved}

\section{Source statement}

Theorem~6.3 of Ref.~\cite{Wolf2012Quantum} concerns an irreducible positive
map
\begin{align}
  T
    &:\MN{D}\longrightarrow\MN{D}.
  \label{eq:wolf_thm6_3_positive_map_cp_scope_positive_map}
\end{align}
It does not assume complete positivity. The exact local source is
\path{Notes/WolfNoteTexSource/ch06_spectral_properties.tex},
lines~604--680.

For each nonzero $X\succeq0$, source Equations~(6.29)--(6.30) define
\begin{align}
  r(X)
    &=\sup\{a\in\mathbb R:T(X)-aX\succeq0\},
  \label{eq:wolf_thm6_3_positive_map_cp_scope_lower_functional}\\
  \widetilde r(X)
    &=\inf\{a\in\mathbb R:T(X)-aX\preceq0\}.
  \label{eq:wolf_thm6_3_positive_map_cp_scope_upper_functional}
\end{align}
The theorem identifies the global lower and upper values, constructs a
strictly positive eigenvector, proves that its ordinary eigenspace is
one-dimensional, compares every positive eigenvalue with the distinguished
value, and identifies that value with the spectral radius
\cite[Theorem~6.3]{Wolf2012Quantum}.

\section{Positive-map scope}

The declaration
\leanid{exists_wolfTheorem63_of_irreducible_positive} in
\doclink{QICLean/Channel/Irreducible/SpectralRadius} proves all four conclusions
for an arbitrary irreducible positive map on a nonzero full matrix algebra. It
produces a density matrix $X>0$ and a scalar $r\geq0$ satisfying
\begin{align}
  T(X)
    &=rX,
  \label{eq:wolf_thm6_3_positive_map_cp_scope_perron_eigenvector}\\
  \varrho(T)
    &=r.
  \label{eq:wolf_thm6_3_positive_map_cp_scope_spectral_radius}
\end{align}
The ordinary complex eigenspace at $r$ has dimension one. Every positive
eigenvalue admitting a nonzero positive semidefinite eigenvector equals $r$.
The global lower value is a maximum and the corrected global upper value is a
minimum, and both are attained at $X$.

The source-form declaration
\leanid{exists_wolfTheorem63_of_irreducible_positive_of_ne_zero} adds the
hypothesis $T\neq0$ and obtains
\begin{align}
  r
    &>0.
  \label{eq:wolf_thm6_3_positive_map_cp_scope_positive_radius}
\end{align}
The structural growth theorem \leanid{growth_posDef_of_irreducible}, the
trace-adjoint theorem \leanid{IsIrreducibleMap.traceAdjointMap}, and the
spectral-radius theorem
\leanid{spectralRadius_eq_of_posDef_eigenvector_of_positive} assume only
positivity. Earlier completely positive declarations remain available as
specializations.

Here ``non-degenerate'' means that the ordinary eigenspace is
one-dimensional, exactly as established by the boundary argument at local
source lines~653--661. No assertion is made about algebraic multiplicity or
the generalized eigenspace.

\section{One-dimensional zero-map boundary}

The printed theorem does not exclude the zero map, but its proof begins at
line~640 with a positive lower value, and item~2 asserts
\begin{align}
  T(X)
    &=rX>0.
  \label{eq:wolf_thm6_3_positive_map_cp_scope_printed_positive_image}
\end{align}
This assertion is false in dimension one. On $\MN{1}$, the zero map is positive
and irreducible under both the source and repository definitions because the
only orthogonal projections are $0$ and $\Id$. Its spectral radius is zero,
and it has no strictly positive eigenvalue.

The assumption-free corrected theorem therefore separates the conclusions
\begin{align}
  X
    &>0,
  \label{eq:wolf_thm6_3_positive_map_cp_scope_corrected_positive_vector}\\
  T(X)
    &=rX,
  \label{eq:wolf_thm6_3_positive_map_cp_scope_corrected_eigenvector}\\
  r
    &\geq0.
  \label{eq:wolf_thm6_3_positive_map_cp_scope_corrected_nonnegative_radius}
\end{align}
The printed conclusion $r>0$ becomes valid after adding $T\neq0$. If $D>1$,
irreducibility already excludes the zero map, but the uniform hypothesis
records the exact missing boundary.

\section{Correction to the upper extremum}

At local source lines~617--619, the printed global definitions and comparison
are
\begin{align}
  r
    &:=\sup_{X\succeq0}r(X),
  \label{eq:wolf_thm6_3_positive_map_cp_scope_printed_lower_global}\\
  \widetilde r
    &:=\sup_{X\succeq0}\widetilde r(X),
  \label{eq:wolf_thm6_3_positive_map_cp_scope_printed_upper_global}\\
  r
    &\geq\widetilde r.
  \label{eq:wolf_thm6_3_positive_map_cp_scope_printed_global_order}
\end{align}
This is not the Collatz--Wielandt pair used by the theorem. The corrected upper
quantity is
\begin{align}
  \widetilde r
    &=\inf_{\substack{X\succeq0\\X\neq0}}\widetilde r(X).
  \label{eq:wolf_thm6_3_positive_map_cp_scope_corrected_upper_global}
\end{align}
Equivalently, the infimum may be taken over density matrices, as in the compact
normalization used by the proof. For an irreducible positive map, the lower
supremum and upper infimum are attained at the same positive-definite
eigenvector and are equal. The reversed inequality at line~619 is part of the
same local source error.

Irreducibility does not imply pointwise equality
$r(X)=\widetilde r(X)$ for every nonzero $X\succeq0$. Consider
$T:\MN{2}\to\MN{2}$ defined by
\begin{align}
  T(Z)
    &=\tr(Z)\Id.
  \label{eq:wolf_thm6_3_positive_map_cp_scope_trace_identity_map}
\end{align}
This map is positive and irreducible: for every nonzero proper projection $P$,
the matrix $T(P)=\Id$ is not supported on the range of $P$. At
\begin{align}
  X
    &=\operatorname{diag}(1,2),
  \label{eq:wolf_thm6_3_positive_map_cp_scope_pointwise_counterexample_matrix}
\end{align}
the lower and upper feasibility conditions give
\begin{align}
  r(X)
    &=\min\{3,3/2\}
     =3/2,
  \label{eq:wolf_thm6_3_positive_map_cp_scope_pointwise_lower_value}\\
  \widetilde r(X)
    &=\max\{3,3/2\}
     =3.
  \label{eq:wolf_thm6_3_positive_map_cp_scope_pointwise_upper_value}
\end{align}
Equality at a positive eigenvector is sufficient for the theorem.

The declarations \leanid{lowerCollatzWielandtValue} and
\leanid{upperCollatzWielandtValue} define the two pointwise functionals in the
extended reals. Their values at the excluded zero matrix are
\begin{align}
  r(0)
    &=+\infty,
  \label{eq:wolf_thm6_3_positive_map_cp_scope_zero_lower_value}\\
  \widetilde r(0)
    &=-\infty.
  \label{eq:wolf_thm6_3_positive_map_cp_scope_zero_upper_value}
\end{align}
If the defining set for the upper functional is empty, its infimum is
$+\infty$. For every nonzero $X\succeq0$, the declaration
\leanid{lowerCollatzWielandtValue\_le\_upperCollatzWielandtValue} proves
\begin{align}
  r(X)
    &\leq\widetilde r(X).
  \label{eq:wolf_thm6_3_positive_map_cp_scope_correct_pointwise_order}
\end{align}
The declaration
\leanid{lower\_and\_upperCollatzWielandtValue\_eq\_of\_eigenvector} proves
equality at every positive eigenvector. Thus the formal result matches the use
at local source lines~649--651 without asserting false pointwise equality for
arbitrary positive semidefinite matrices.

\section{Source-faithful proof route}

The formal proof follows the order of the source argument
\cite[Theorem~6.3]{Wolf2012Quantum}. Lower and upper feasibility are first
defined for every real scalar as in
(\ref{eq:wolf_thm6_3_positive_map_cp_scope_lower_functional}) and
(\ref{eq:wolf_thm6_3_positive_map_cp_scope_upper_functional}). After
normalization to density matrices, zero is lower feasible and negative lower
candidates cannot increase the supremum. The maximizing search can therefore
be restricted to a compact interval. Positivity and the trace inequality make
all upper candidates nonnegative. These observations produce a lower maximizer
without changing either source extremum.

Source Equation~(6.31) of Ref.~\cite{Wolf2012Quantum}, together with the positive-map growth theorem, makes
the maximizer positive definite and forces its lower residual to vanish:
\begin{align}
  T(X)-rX
    &=0.
  \label{eq:wolf_thm6_3_positive_map_cp_scope_vanishing_lower_residual}
\end{align}
Source Equation~(6.32) of Ref.~\cite{Wolf2012Quantum} is then applied after decomposing an arbitrary matrix
into Hermitian and skew-Hermitian parts. This proves that the ordinary
eigenspace is one-dimensional.

If $r=0$, positivity and $X>0$ imply that a map satisfying $T(X)=0$ is the
zero map, after which the remaining conclusions are immediate. In the branch
$r>0$, the complementary-projection argument at local source lines~604--606
shows that $T^*$ is irreducible. A positive-definite Perron eigenvector for
$T^*$ supplies source Equation~(6.33) of Ref.~\cite{Wolf2012Quantum}. Pairing with this vector proves
uniqueness of positive eigenvalues and, after pairing with an upper residual,
the corrected upper minimum.

Finally, conjugation by $X^{1/2}$ and rescaling by $r^{-1}$ produce the
positive unital map used by Wolf. Proposition~6.1 of
Ref.~\cite{Wolf2012Quantum} identifies its spectral radius as one, which gives
(\ref{eq:wolf_thm6_3_positive_map_cp_scope_spectral_radius}).

\section{Verdict}

\textbf{Verdict: resolved, with two source corrections.} The former
complete-positivity restriction has been removed: Theorem~6.3 is formalized
for arbitrary positive maps. At the one-dimensional zero-map boundary, the
assumption-free theorem requires $r\geq0$; the printed conclusion $r>0$
requires the additional hypothesis $T\neq0$. The second global extremum at
local source lines~617--619 is an infimum rather than a supremum, and the
printed comparison is correspondingly invalid. The pointwise functionals use
extended values, their zero and empty-set boundaries are explicit, and their
equality is asserted only at positive eigenvectors.

\bibliographystyle{alpha}
\bibliography{references}

\end{document}
